(1) 能量守恒给 $mv^2/2=mgy$，因此 $v=\sqrt{2gy}$。线元 $ds=\sqrt{dx^2+dy^2}=\sqrt{1+y'^2}\,dx$，行进时间 $dt=ds/v$，于是 $T[y]=\int_{x_0}^{x_1}\sqrt{1+y'^2}\,(2gy)^{-1/2}dx$。端点附近 $y=0$ 的积分需按可积的反常积分理解。
(2) 被积函数不显含 $x$，Beltrami恒等式 $F-y'F_{y'}=\text{常数}$ 给 $[y(1+y'^2)]^{-1/2}=(2a)^{-1/2}$，即 $y(1+y'^2)=2a$。取 $y=a(1-\cos\theta)$，则 $dx/d\theta=a(1-\cos\theta)$、$dy/d\theta=a\sin\theta$，积分并固定起点得 $x=a(\theta-\sin\theta)$、$y=a(1-\cos\theta)$，满足首次积分。
(3) 将摆线倒置，使最低点位于参数 $\theta=\pi$，令高度相对最低点为 $h=a(1+\cos\theta)$。沿曲线的弧长微分 $ds/d\theta=2a\sin(\theta/2)$（取 $0<\theta<2\pi$），从最低点起的有符号弧长可写 $s=-4a\cos(\theta/2)$。由 $\cos\theta=2\cos^2(\theta/2)-1$ 得 $h=2a\cos^2(\theta/2)=s^2/(8a)$，故 $U=mgs^2/(8a)$。
(4) 弧长方向的动力学为 $m\ddot s=-dU/ds=-mgs/(4a)$，所以角频率 $\omega=\sqrt{g/(4a)}$，与振幅无关。从任意非零振幅静止释放到最低点经历四分之一周期 $T_{\downarrow}=\pi/(2\omega)=\pi\sqrt{a/g}$，这是理想无摩擦摆线槽的等时性；它不等同于普通圆弧单摆在大幅度时的运动。
